

Deep learning techniques have greatly improved medical image segmentation by automatically extracting specific tissue or substance location information, which facilitates accurate disease diagnosis and assessment. However, most deep learning approaches for segmentation require fully or partially labeled training datasets, which can be time-consuming and expensive to annotate. To address this issue, recent research has focused on developing segmentation frameworks that require little or no segmentation labels.

To meet this need, many researchers have devoted their efforts to Weakly-Supervised Semantic Segmentation (WSSS)\cite{wolleb2022diffusion}, which utilizes weak supervision, such as image-level classification labels. Recent WSSS methods can be broadly categorized into two types \cite{chan2021comprehensive}: Class-Activation-Mapping-based (CAM-based) \cite{selvaraju2017grad,wang2020score,jiang2021layercam,wang2020self,lee2021lfi,xie2022c2am}, and Multiple-Instance-Learning-based (MIL-based) \cite{qian2022transformer} methods.

The literature has not adequately addressed the issue of low-resolution Class-Activation Maps (CAMs), especially for medical images. Some existing methods, such as dilated residual networks \cite{yu2017dilated} and U-Net segmentation architecture \cite{belharbi2022f,englebert2022poly,tagaris2019high}, have attempted to tackle this issue, but still require many upsampling operations, which the results become blurry. Meanwhile, LayerCAM \cite{jiang2021layercam} has proposed a hierarchical solution that extracts activation maps from multiple convolution layers using Grad-CAM\cite{selvaraju2017grad} and aggregates them with equal weights. Although this approach successfully enhances the resolution of the segmentation mask, it lacks flexibility and may not be optimal.

In this paper, we propose an Attentive Multiple-Exit CAM (AME-CAM) for brain tumor segmentation in magnetic resonance imaging (MRI). Different from recent CAM methods, AME-CAM uses a classification model with multiple-exit training strategy applied to optimize the internal outputs. Activation maps from the outputs of internal classifiers, which have different resolutions, are then aggregated using an attention model. The model learns the pixel-wise weighted sum of the activation maps by a novel contrastive learning method.

Our proposed method has the following contributions:
\begin{itemize}
\item To tackle the issues in existing CAMs, we propose to use multiple-exit classification networks to accurately capture all the internal activation maps of different resolutions.
\item We propose an attentive feature aggregation to learn the pixel-wise weighted sum of the internal activation maps.
\item We demonstrate the superiority of AME-CAM over state-of-the-art CAM methods in extracting segmentation results from classification networks on the 2021 Brain Tumor Segmentation Challenge (BraTS 2021) \cite{menze2014multimodal,bakas2017advancing,bakas2018identifying}.
\item For reproducibility, we have released our code at \\ \href{https://github.com/windstormer/AME-CAM}{https://github.com/windstormer/AME-CAM}
\end{itemize}

Overall, our proposed method can help overcome the challenges of expensive and time-consuming segmentation labeling in medical imaging, and has the potential to improve the accuracy of disease diagnosis and assessment.